% --- Archon physics formalization source begin ---
% archon:physics
% archon:covers PhyXMiniProblems/problem_phyx_mini_0474.lean
% archon:source-report reports/phyx_mini/problem_phyx_mini_0474.source.json

\chapter{Physics problem phyx\_mini\_0474}
\label{ch:PhyXMiniProblems_problem_phyx_mini_0474}

\paragraph{Problem source.}
\#\# Physical scenario

A Carnot engine takes 2000 J of heat from a reservoir at 500 K, does some work, and discards some heat to a reservoir at 350 K.

\#\# Question

What is its efficiency?

\#\# Answer choices

- A: A: 36\%

- B: B: 30\%

- C: C: 28\%

- D: D: 45\%

\#\# Figure caption (auxiliary; use the image as primary evidence)

The diagram depicts a heat engine process involving heat transfer and work done. Here's a breakdown of the content:

- **Arrows and Labels**:
  - A vertical arrow pointing downward from \textbackslash{}( T\_H \textbackslash{}) to \textbackslash{}( T\_C \textbackslash{}), representing heat transfer.
  - A horizontal arrow pointing to the right representing work output.

- **Textual Information**:
  - Top of the diagram: \textbackslash{}( T\_H = 500 \textbackslash{}, \textbackslash{}text\{K\} \textbackslash{}) indicating the high temperature reservoir.
  - Below \textbackslash{}( T\_H \textbackslash{}): \textbackslash{}( Q\_H = 2000 \textbackslash{}, \textbackslash{}text\{J\} \textbackslash{}) indicating the amount of heat absorbed from the high temperature reservoir.
  - Below the horizontal arrow: \textbackslash{}( W = ? \textbackslash{}) and \textbackslash{}( e = ? \textbackslash{}), representing the unknown work done by the engine and the efficiency.
  - Bottom of the diagram: \textbackslash{}( T\_C = 350 \textbackslash{}, \textbackslash{}text\{K\} \textbackslash{}) indicating the low temperature reservoir.
  - Beside the downward arrow near \textbackslash{}( T\_C \textbackslash{}): \textbackslash{}( Q\_C = ? \textbackslash{}), representing the unknown amount of heat expelled to the low temperature reservoir.

This diagram visually represents the key components and relationships in a heat engine cycle.

\#\# Dataset metadata

Laws of Thermodynamics; Physical Model Grounding Reasoning, Multi-Formula Reasoning

\paragraph{Recorded answer/context.}
B: B: 30\%

\paragraph{Figure/image path.}
/root/proposal\_for\_physic/science-mango/phyx\_mini\_run/phyx\_data/test\_image/474.png

\paragraph{Formalization target.}
create a compiling Lean file with sorry bodies at `PhyXMiniProblems/problem\_phyx\_mini\_0474.lean`. The Lean declarations must preserve the physical quantities, dimensions or dimensional roles, figure labels, governing-law hypotheses, and final relation expressed by this problem.
Use Mathlib/Physlib names found through LeanExplore where available. If a physics API is missing, introduce faithful local abstractions rather than scalar placeholder aliases.

\begin{theorem}[Physics formalization target]
\label{thm:physics:phyx_mini_0474:target}
\lean{PhyXMiniProblems.ProblemPhyXMini0474.carnotEfficiency_eq_recordedAnswerB}
\uses{def:physics:phyx-mini-0474:phyxminiproblems-problemphyxmini0474-carnotheatenginesetup, def:physics:phyx-mini-0474:phyxminiproblems-problemphyxmini0474-matchesproblemstatement, def:physics:phyx-mini-0474:phyxminiproblems-problemphyxmini0474-matchesprimaryheatenginefigure, def:physics:phyx-mini-0474:phyxminiproblems-problemphyxmini0474-hasphysicalcarnotengineparameters, def:physics:phyx-mini-0474:phyxminiproblems-problemphyxmini0474-satisfiescarnotheatenginelaws, def:physics:phyx-mini-0474:phyxminiproblems-problemphyxmini0474-recordedanswerchoice}
The assigned autoformalize agent should translate this physics problem into Lean declarations in the covered file, with theorem and lemma proof bodies written as `by sorry`.
\end{theorem}
\begin{proof}
This is an autoformalization task, not a proof task. Produce statements that can later be proved by the physics prover without weakening the physical model.
\end{proof}

% --- Archon declaration topology begin ---
\section{Declaration topology}
\begin{definition}[Temperature Quantity]
  \label{def:physics:phyx-mini-0474:phyxminiproblems-problemphyxmini0474-temperaturequantity}
  \lean{PhyXMiniProblems.ProblemPhyXMini0474.TemperatureQuantity}
  Absolute thermodynamic temperature
\end{definition}
\begin{proof}
  This is the declared model definition of Temperature Quantity.
\end{proof}
\begin{definition}[Energy Quantity]
  \label{def:physics:phyx-mini-0474:phyxminiproblems-problemphyxmini0474-energyquantity}
  \lean{PhyXMiniProblems.ProblemPhyXMini0474.EnergyQuantity}
  A physical heat or work quantity, carrying the dimension of energy
\end{definition}
\begin{proof}
  This is the declared model definition of Energy Quantity.
\end{proof}
\begin{definition}[energy In Joules]
  \label{def:physics:phyx-mini-0474:phyxminiproblems-problemphyxmini0474-energyinjoules}
  \lean{PhyXMiniProblems.ProblemPhyXMini0474.energyInJoules}
  \uses{def:physics:phyx-mini-0474:phyxminiproblems-problemphyxmini0474-energyquantity}
  Read a physical energy quantity in joules
\end{definition}
\begin{proof}
  This is the declared model definition of energy In Joules.
\end{proof}
\begin{definition}[temperature In Kelvins]
  \label{def:physics:phyx-mini-0474:phyxminiproblems-problemphyxmini0474-temperatureinkelvins}
  \lean{PhyXMiniProblems.ProblemPhyXMini0474.temperatureInKelvins}
  \uses{def:physics:phyx-mini-0474:phyxminiproblems-problemphyxmini0474-temperaturequantity}
  Read an absolute temperature in kelvins from its stated storage unit
\end{definition}
\begin{proof}
  This is the declared model definition of temperature In Kelvins.
\end{proof}
\begin{definition}[Thermodynamic Device Role]
  \label{def:physics:phyx-mini-0474:phyxminiproblems-problemphyxmini0474-thermodynamicdevicerole}
  \lean{PhyXMiniProblems.ProblemPhyXMini0474.ThermodynamicDeviceRole}
  The thermodynamic role of the device in the problem
\end{definition}
\begin{proof}
  This is the declared model definition of Thermodynamic Device Role.
\end{proof}
\begin{definition}[Heat Engine Model]
  \label{def:physics:phyx-mini-0474:phyxminiproblems-problemphyxmini0474-heatenginemodel}
  \lean{PhyXMiniProblems.ProblemPhyXMini0474.HeatEngineModel}
  The reversible ideal engine model specified by the problem
\end{definition}
\begin{proof}
  This is the declared model definition of Heat Engine Model.
\end{proof}
\begin{definition}[Reservoir]
  \label{def:physics:phyx-mini-0474:phyxminiproblems-problemphyxmini0474-reservoir}
  \lean{PhyXMiniProblems.ProblemPhyXMini0474.Reservoir}
  The two reservoirs with which the engine exchanges heat
\end{definition}
\begin{proof}
  This is the declared model definition of Reservoir.
\end{proof}
\begin{definition}[Diagram Node]
  \label{def:physics:phyx-mini-0474:phyxminiproblems-problemphyxmini0474-diagramnode}
  \lean{PhyXMiniProblems.ProblemPhyXMini0474.DiagramNode}
  Locations joined by the arrows in the primary diagram
\end{definition}
\begin{proof}
  This is the declared model definition of Diagram Node.
\end{proof}
\begin{definition}[Diagram Arrow]
  \label{def:physics:phyx-mini-0474:phyxminiproblems-problemphyxmini0474-diagramarrow}
  \lean{PhyXMiniProblems.ProblemPhyXMini0474.DiagramArrow}
  Directed energy-transfer arrows shown in the primary diagram
\end{definition}
\begin{proof}
  This is the declared model definition of Diagram Arrow.
\end{proof}
\begin{definition}[Arrow Orientation]
  \label{def:physics:phyx-mini-0474:phyxminiproblems-problemphyxmini0474-arroworientation}
  \lean{PhyXMiniProblems.ProblemPhyXMini0474.ArrowOrientation}
  Geometric orientation of an arrow in the primary diagram
\end{definition}
\begin{proof}
  This is the declared model definition of Arrow Orientation.
\end{proof}
\begin{definition}[Figure Label]
  \label{def:physics:phyx-mini-0474:phyxminiproblems-problemphyxmini0474-figurelabel}
  \lean{PhyXMiniProblems.ProblemPhyXMini0474.FigureLabel}
  Textual quantity labels printed in the primary diagram
\end{definition}
\begin{proof}
  This is the declared model definition of Figure Label.
\end{proof}
\begin{definition}[Heat Engine Diagram]
  \label{def:physics:phyx-mini-0474:phyxminiproblems-problemphyxmini0474-heatenginediagram}
  \lean{PhyXMiniProblems.ProblemPhyXMini0474.HeatEngineDiagram}
  \uses{def:physics:phyx-mini-0474:phyxminiproblems-problemphyxmini0474-diagramnode, def:physics:phyx-mini-0474:phyxminiproblems-problemphyxmini0474-diagramarrow, def:physics:phyx-mini-0474:phyxminiproblems-problemphyxmini0474-arroworientation, def:physics:phyx-mini-0474:phyxminiproblems-problemphyxmini0474-figurelabel}
  Raw label and arrow information transcribed from the primary bitmap
\end{definition}
\begin{proof}
  This is the declared model definition of Heat Engine Diagram.
\end{proof}
\begin{definition}[Carnot Heat Engine Setup]
  \label{def:physics:phyx-mini-0474:phyxminiproblems-problemphyxmini0474-carnotheatenginesetup}
  \lean{PhyXMiniProblems.ProblemPhyXMini0474.CarnotHeatEngineSetup}
  \uses{def:physics:phyx-mini-0474:phyxminiproblems-problemphyxmini0474-temperaturequantity, def:physics:phyx-mini-0474:phyxminiproblems-problemphyxmini0474-energyquantity, def:physics:phyx-mini-0474:phyxminiproblems-problemphyxmini0474-thermodynamicdevicerole, def:physics:phyx-mini-0474:phyxminiproblems-problemphyxmini0474-heatenginemodel, def:physics:phyx-mini-0474:phyxminiproblems-problemphyxmini0474-reservoir, def:physics:phyx-mini-0474:phyxminiproblems-problemphyxmini0474-heatenginediagram}
  The physical Carnot-engine setup. Reservoir heat transfers and net work are stored as positive energy magnitudes; the efficiency is dimensionless
\end{definition}
\begin{proof}
  This is the declared model definition of Carnot Heat Engine Setup.
\end{proof}
\begin{definition}[Matches Problem Statement]
  \label{def:physics:phyx-mini-0474:phyxminiproblems-problemphyxmini0474-matchesproblemstatement}
  \lean{PhyXMiniProblems.ProblemPhyXMini0474.MatchesProblemStatement}
  \uses{def:physics:phyx-mini-0474:phyxminiproblems-problemphyxmini0474-carnotheatenginesetup}
  The prose identifies the device as a Carnot heat engine
\end{definition}
\begin{proof}
  This is the declared model definition of Matches Problem Statement.
\end{proof}
\begin{definition}[Matches Primary Heat Engine Figure]
  \label{def:physics:phyx-mini-0474:phyxminiproblems-problemphyxmini0474-matchesprimaryheatenginefigure}
  \lean{PhyXMiniProblems.ProblemPhyXMini0474.MatchesPrimaryHeatEngineFigure}
  \uses{def:physics:phyx-mini-0474:phyxminiproblems-problemphyxmini0474-energyinjoules, def:physics:phyx-mini-0474:phyxminiproblems-problemphyxmini0474-temperatureinkelvins, def:physics:phyx-mini-0474:phyxminiproblems-problemphyxmini0474-figurelabel, def:physics:phyx-mini-0474:phyxminiproblems-problemphyxmini0474-carnotheatenginesetup}
  Figure-derived data and geometry. The image supplies only the hot input heat and the two reservoir temperatures; rejected heat, work, and efficiency remain explicitly marked unknown
\end{definition}
\begin{proof}
  This is the declared model definition of Matches Primary Heat Engine Figure.
\end{proof}
\begin{definition}[Has Physical Carnot Engine Parameters]
  \label{def:physics:phyx-mini-0474:phyxminiproblems-problemphyxmini0474-hasphysicalcarnotengineparameters}
  \lean{PhyXMiniProblems.ProblemPhyXMini0474.HasPhysicalCarnotEngineParameters}
  \uses{def:physics:phyx-mini-0474:phyxminiproblems-problemphyxmini0474-energyinjoules, def:physics:phyx-mini-0474:phyxminiproblems-problemphyxmini0474-temperatureinkelvins, def:physics:phyx-mini-0474:phyxminiproblems-problemphyxmini0474-carnotheatenginesetup}
  Positivity and ordering conditions for a physically operating engine
\end{definition}
\begin{proof}
  This is the declared model definition of Has Physical Carnot Engine Parameters.
\end{proof}
\begin{definition}[Satisfies Carnot Heat Engine Laws]
  \label{def:physics:phyx-mini-0474:phyxminiproblems-problemphyxmini0474-satisfiescarnotheatenginelaws}
  \lean{PhyXMiniProblems.ProblemPhyXMini0474.SatisfiesCarnotHeatEngineLaws}
  \uses{def:physics:phyx-mini-0474:phyxminiproblems-problemphyxmini0474-energyinjoules, def:physics:phyx-mini-0474:phyxminiproblems-problemphyxmini0474-temperatureinkelvins, def:physics:phyx-mini-0474:phyxminiproblems-problemphyxmini0474-carnotheatenginesetup}
  Governing relations for one complete reversible Carnot cycle. Heat magnitudes are positive: the first law gives W = Q\_H - Q\_C, efficiency is W / Q\_H, and reversibility gives Q\_C / Q\_H = T\_C / T\_H. None of these fields fixes the requested numerical efficiency
\end{definition}
\begin{proof}
  This is the declared model definition of Satisfies Carnot Heat Engine Laws.
\end{proof}
\begin{lemma}[rejected Heat in Joules eq 1400]
  \label{lem:physics:phyx-mini-0474:phyxminiproblems-problemphyxmini0474-rejectedheat-injoules-eq-1400}
  \lean{PhyXMiniProblems.ProblemPhyXMini0474.rejectedHeat_inJoules_eq_1400}
  \uses{def:physics:phyx-mini-0474:phyxminiproblems-problemphyxmini0474-energyinjoules, def:physics:phyx-mini-0474:phyxminiproblems-problemphyxmini0474-carnotheatenginesetup, def:physics:phyx-mini-0474:phyxminiproblems-problemphyxmini0474-matchesproblemstatement, def:physics:phyx-mini-0474:phyxminiproblems-problemphyxmini0474-matchesprimaryheatenginefigure, def:physics:phyx-mini-0474:phyxminiproblems-problemphyxmini0474-hasphysicalcarnotengineparameters, def:physics:phyx-mini-0474:phyxminiproblems-problemphyxmini0474-satisfiescarnotheatenginelaws}
  The reversible heat ratio determines 1400 J of rejected heat
\end{lemma}
\begin{proof}
  The conclusion follows from the governing relations and figure readouts named in the statement of rejected Heat in Joules eq 1400.
\end{proof}
\begin{lemma}[net Work Output in Joules eq 600]
  \label{lem:physics:phyx-mini-0474:phyxminiproblems-problemphyxmini0474-networkoutput-injoules-eq-600}
  \lean{PhyXMiniProblems.ProblemPhyXMini0474.netWorkOutput_inJoules_eq_600}
  \uses{def:physics:phyx-mini-0474:phyxminiproblems-problemphyxmini0474-energyinjoules, def:physics:phyx-mini-0474:phyxminiproblems-problemphyxmini0474-carnotheatenginesetup, def:physics:phyx-mini-0474:phyxminiproblems-problemphyxmini0474-matchesproblemstatement, def:physics:phyx-mini-0474:phyxminiproblems-problemphyxmini0474-matchesprimaryheatenginefigure, def:physics:phyx-mini-0474:phyxminiproblems-problemphyxmini0474-hasphysicalcarnotengineparameters, def:physics:phyx-mini-0474:phyxminiproblems-problemphyxmini0474-satisfiescarnotheatenginelaws}
  The first law then determines 600 J of net work output
\end{lemma}
\begin{proof}
  The conclusion follows from the governing relations and figure readouts named in the statement of net Work Output in Joules eq 600.
\end{proof}
\begin{definition}[Answer Choice]
  \label{def:physics:phyx-mini-0474:phyxminiproblems-problemphyxmini0474-answerchoice}
  \lean{PhyXMiniProblems.ProblemPhyXMini0474.AnswerChoice}
  Labels attached to the four displayed multiple-choice answers
\end{definition}
\begin{proof}
  This is the declared model definition of Answer Choice.
\end{proof}
\begin{definition}[efficiency Percent]
  \label{def:physics:phyx-mini-0474:phyxminiproblems-problemphyxmini0474-answerchoice-efficiencypercent}
  \lean{PhyXMiniProblems.ProblemPhyXMini0474.AnswerChoice.efficiencyPercent}
  \uses{def:physics:phyx-mini-0474:phyxminiproblems-problemphyxmini0474-answerchoice}
  Efficiency percentage printed beside each answer label
\end{definition}
\begin{proof}
  This is the declared model definition of efficiency Percent.
\end{proof}
\begin{definition}[recorded Answer Choice]
  \label{def:physics:phyx-mini-0474:phyxminiproblems-problemphyxmini0474-recordedanswerchoice}
  \lean{PhyXMiniProblems.ProblemPhyXMini0474.recordedAnswerChoice}
  \uses{def:physics:phyx-mini-0474:phyxminiproblems-problemphyxmini0474-answerchoice}
  The answer label recorded in the dataset metadata
\end{definition}
\begin{proof}
  This is the declared model definition of recorded Answer Choice.
\end{proof}
% --- Archon declaration topology end ---
% --- Archon physics formalization source end ---
